\documentclass[aspectratio=169]{beamer}
\usetheme{Madrid}
\usecolortheme{beaver}
\usepackage[utf8]{inputenc}
\usepackage[T1]{fontenc}
\setbeamertemplate{navigation symbols}{}
\setbeamercolor{title}{fg=white,bg=cyan!60!black}
\setbeamercolor{frametitle}{fg=white,bg=cyan!50!black}
\setbeamercolor{block title}{bg=cyan!15,fg=black}
\setbeamercolor{block body}{bg=cyan!5,fg=black}

\title{Block Arcade 1 — Sprites \& the Stage}
\subtitle{Every sprite lives at an (x, y) — coordinates in blocks}
\author{ByteForge --- Grade 7 --- Game Maker}
\date{}

\begin{document}

\begin{frame}
    \titlepage
\end{frame}

\begin{frame}{Today's Plan}
    \begin{itemize}
        \item Meet Your Sprite
        \item Two Sprites, Two Corners
        \item[$\bigstar$] BOSS --- Place Your Cast
    \end{itemize}
\end{frame}

\begin{frame}{A sprite is just a shape at a position}
    \Large
    A "sprite" is any character or object on the stage — under the hood, it's a position (x, y) plus something drawn there.
\end{frame}


\begin{frame}{Naming positions makes them easier to reason about}
    \Large
    "Top-left" and "bottom-right" are shorthand for specific (x,y) values — translating a description into coordinates is a skill you'll use constantly.
\end{frame}


\begin{frame}{BOSS — three sprites, three positions}
    \Large
    Same skill, just repeated — real Scratch projects often start with exactly this.
\end{frame}

\begin{frame}{Real-World Example}
    \Large
    Open Scratch's own stage and drag a sprite around — the x/y readout in the corner is literally the same two numbers this chapter is teaching.
\end{frame}

\begin{frame}{Watch}
    \Large
    Suggested video topic: \emph{Scratch coordinates explained — x and y for beginners}

    \vspace{1em}
    \normalsize
    Search this on YouTube together, or swap in a specific video you've
    already vetted --- do not project an unverified link cold in class.
\end{frame}

\begin{frame}{Let Them Discover}
    \Large
    On a printed 480x260 grid, mark where you'd place a sprite for "top-left corner", "dead center", and "bottom-right corner" before touching the platform — compare answers as a class.
\end{frame}

\begin{frame}{Your Turn}
    \Large
    Open the student map, find \textbf{Block Arcade 1 — Sprites \& the Stage}, and work through it
    hands-on --- everything you just saw, they now do themselves.
\end{frame}

\end{document}
